Man rechnet jeweils beide Grössen in dieselbe Einheit um.
\begin{abcliste}
\abc $\aLiO \;\result{\neq}\; \aReO$: eine Masse lässt sich nicht mit einer Kraft vergleichen.
\abc $\bLiO = \bLi \;\result{=}\; \bReO = \bRe$
\abc $\cLiO \;\result{=}\; \cReO$, denn $\si{\newton\second\squared\per\meter} = \si{\kilo\gram\meter\per\second\squared}\cdot\si{\second\squared\per\meter} = \si{\kilo\gram}$
\abc $\dLiO = \dLsP{-9}{2} \;\result{>}\; \dReO$
\abc $\eReO = \eRsP{0}{2}$, also $\eLiO \;\result{=}\; \eReO$
\abc $\fAO\cdot\fBO = \fA\cdot\fB = \fLmP{0}{2} \;\result{<}\; \fReO$
\end{abcliste}
